% T02 source: Базовая модель.tex lines 90--183
\subsection{Level M0: Ideally Divisible Work}

\begin{lemma}[Baseline Time without AI]\label{lem:baseline}
When the tasks are performed sequentially by a developer without using AI, the total time is
\begin{equation}
  T_h = \sum_{i=1}^{N} Z_i X.
\end{equation}
\end{lemma}

\begin{assumption}[Idealized AI-Assisted Time Model (Level M0)]\label{assum:ideal}
In the core analytical model, the work is assumed to be ideally divisible, with parallelism distributed uniformly. Define the total AI-assisted workload as
\[
  W := X \sum_{i=1}^{N} Z_i k_i.
\]
When AI is used with parallelism factor $P$, the total task-completion time is taken to be
\begin{equation}
  T_{ai}^{(0)} = \frac{W}{P} = \frac{1}{P} \sum_{i=1}^{N} Z_i k_i X.
\end{equation}
All subsequent theoretical results are based on this assumption unless another level of the model hierarchy is explicitly specified (Remark~\ref{rem:hierarchy}).
\end{assumption}

\begin{observation}[Speedup Criterion and Factor at Level M0]\label{obs:ideal}
Define the weighted-average normalized duration, with task sizes $Z_i$ as weights, as
\[
  \bar{k}_w := \frac{\sum_{i=1}^{N} Z_i k_i}{\sum_{i=1}^{N} Z_i}.
\]
Substitution from Lemma~\ref{lem:baseline} and Assumption~\ref{assum:ideal} directly yields the chain of equivalences
\begin{equation}
  T_{ai}^{(0)} < T_h
  \quad\Longleftrightarrow\quad
  \sum_{i=1}^{N} Z_i k_i < P \sum_{i=1}^{N} Z_i
  \quad\Longleftrightarrow\quad
  P > \bar{k}_w.
\end{equation}
This is not an independent theoretical result, but an algebraic restatement of the definition of time at Level M0.

The speedup factor is
\begin{equation}
  S_E = \frac{T_h}{T_{ai}^{(0)}} = P \cdot \frac{\sum_{i=1}^{N} Z_i}{\sum_{i=1}^{N} Z_i k_i} = \frac{P}{\bar{k}_w}.
\end{equation}
In particular, if $k_i \le 1$ for every $i$ and $P > 1$, then $S_E \ge P > 1$; that is, speedup is guaranteed.

\medskip
Furthermore, $\min_{j} k_j \le \bar{k}_w \le \max_{j} k_j$ implies the two-sided bound
\[
  \frac{P}{\max_{i} k_i} \;\le\; S_E \;\le\; \frac{P}{\min_{i} k_i},
\]
which provides a quick estimate of the speedup range without weighting by $Z_i$.
\end{observation}

\begin{remark}[Speedup Factorization and Baseline Choice]\label{rem:factorization}
The formula in Observation~\ref{obs:ideal} admits the factorization
\[
  S_E = \underbrace{P\vphantom{\frac{1}{\bar{k}_w}}}_{\text{organizational factor}} \cdot \underbrace{\frac{1}{\bar{k}_w}}_{\text{tool factor}}.
\]
The first factor captures the effect of organizing work in parallel, whereas the second captures the effect of AI on task effort. Sequential work by a single developer ($P_h = 1$) is deliberately chosen as the no-AI baseline $T_h$. The model therefore addresses what a single developer gains by transitioning to AI agents, rather than comparing AI with a team of $P$ people. If $k_i \equiv 1$, then $S_E = P$: this speedup is due entirely to process organization and can, in principle, also be attained by a team of $P$ people. The central assumption is that the $P$ parallel streams are provided by AI agents supervised by the same developer, without involving additional specialists. Supervision overhead is accounted for at Levels M3--M4: the function $C(P)$ represents cross-stream integration overhead (Remark~\ref{rem:coordination}), whereas the term $\gamma(P)\,H$ represents the human resource and its associated ceiling $S_E^{\mathrm{ach}} \le 1/\bar{h}_w$ (Remark~\ref{rem:human}, Proposition~\ref{prop:amdahl}). For comparison with a team of $q$ developers, the baseline time should be replaced, under the idealized approximation, by $T_h(q) = T_h/q$; the condition under which the AI-assisted process is superior then becomes $P/\bar{k}_w > q$.

\medskip
The factorization is valid \emph{for each fixed configuration} $\sigma = (P, r)$, but its factors cannot be varied independently. In practice, changing $P$ is often accompanied by a change in operating mode $r$, for example, a transition from interactive work to autonomous agents. The change in mode changes the entire vector $\{k_i^{(r)}\}$ and hence $\bar{k}_w$. Consequently, $\bar{k}_w$ must not be estimated in one mode and combined with the parallelism of another; for example, one must not measure $k_i$ during interactive work and substitute $P = 4$ for autonomous agents. Configurations must be compared as a whole: $S_E(\sigma_1)$ and $S_E(\sigma_2)$.
\end{remark}

\begin{remark}[Critical Parallelism Threshold for $k_i > 1$]\label{rem:threshold}
The identity in Observation~\ref{obs:ideal},
\[
  \sum_{i=1}^{N} Z_i k_i < P \sum_{i=1}^{N} Z_i,
\]
implies that a speedup $S_E > 1$ is achieved if and only if
\[
  P > \bar{k}_w,
\]
where $\bar{k}_w$ is the weighted-average normalized duration introduced in Observation~\ref{obs:ideal}.

We emphasize that the condition $P > \bar{k}_w$ is necessary and sufficient only at Level M0. At Levels M1--M2, it remains necessary but is no longer sufficient: the achievable time is also bounded below by the longest task ($M_N$, Remark~\ref{rem:indivisible}) and by the critical path ($L$, Remark~\ref{rem:graph}); Theorem~\ref{thm:sufficient} provides a substantive sufficient result for these levels. At Level M3, the necessary condition based on the aggregate workload of the streams is $C(P)\bar a_w+\bar h_w<P$ (Remark~\ref{rem:coordination}). At Level M4, all three inequalities specified by $B_4<T_h$ are simultaneously necessary: $W_4/P<T_h$, $L_4<T_h$, and $\gamma(P)H<T_h$. They are not sufficient because there exist instances in which $T^{*}>B_4$ due to local blocking and the queue for developer attention.

\begin{itemize}
  \item If $k_i \equiv k > 1$ for every $i$ (AI uniformly slows execution), the critical threshold simplifies to
  \[
    P_{\text{crit}} = k.
  \]
  Thus, parallelism must \textbf{strictly exceed} the slowdown factor; otherwise, using AI provides no elapsed-time advantage.

  \item If the coefficients $k_i$ are heterogeneous, the threshold is determined by the weighted average: tasks with larger task sizes contribute more to $\bar{k}_w$. Thus, even when some tasks have $k_i \gg 1$, speedup is possible if their contribution to the total task size is small and parallelism is sufficient to compensate for it.

  \item In practical applications, $\bar{k}_w$ characterizes the duration of the selected AI-assisted configuration for a given project. The estimation procedure is presented in Remark~\ref{rem:estimation}. If $\bar{k}_w$ is no smaller than the configured parallelism $P$, this configuration provides no elapsed-time gain. In that case, one should
  \begin{enumerate}
    \item reconsider the set of tasks selected for automation, beginning with those having the largest $k_i$;
    \item increase the permissible process parallelism, for example through asynchronous execution;
    \item reduce $k_i$ by changing how AI is used, including prompt formulation, context composition, and postprocessing.
  \end{enumerate}
\end{itemize}

Thus, rather than providing a binary assessment of the usefulness of AI, the model specifies a \textbf{quantitative decision criterion} that relates durations in the selected mode ($k_i^{(r)}$), project structure ($Z_i$), and configured parallelism ($P$). The criterion is applied to the configuration $\sigma = (P, r)$ as a whole: the values of $\bar{k}_w$ and $P$ must refer to the same operating mode (Remark~\ref{rem:factorization}). The choice between fully manual and AI-assisted execution of individual tasks is outside the scope of the formulation; once that choice has been made, the model schedules only the included AI-assisted task set.
\end{remark}
